\begin{textsample}{POS Dim 3 – 2021\_c5 – Score 2.00 – 2021\_c5/enade\_2021\_ci\_soc\_edcb3715cc.txt}  \label{ex:f3_pos_085}
CAPÍTULO 5 DISTRIBUIÇÃO DOS CONCEITOS 5.1 PANORAMA NACIONAL DA DISTRIBUIÇÃO DOS CONCEITOS O Conceito Enade é um indicador de qualidade que avalia os cursos por intermédio dos desempenhos dos estudantes a partir dos resultados do Enade.
Os valores possíveis vão de 1 ( pior situação ) a 5 ( melhor situação ).
Os cursos com apenas um ou sem nenhum concluinte participante não obtêm o Conceito Enade, ficando Sem Conceito ( SC ).
Foram avaliados em Ciências Sociais 164 cursos, sendo 93 ( 56,7 \% ) de Licenciatura e 71 ( 43,3 \% ) de Bacharelado.
Enfatiza -se, mais uma vez, que a diferença entre os cursos tabulados neste capítulo e no Capítulo 2 corresponde aos cursos que não foram avaliados, ou seja, àqueles, em princípio, sem alunos concluintes que tivessem sido inscritos para a prova.
Na Tabela 5.1a, são apresentadas a quantidade e a distribuição de cursos de Licenciatura de Ciências Sociais participantes do Enade/2021, por faixa de conceito e Grande Região.
Observando -se os dados da Tabela 5.1a, nota -se que, dos 93 cursos participantes, 30 ( 32,3 \% ) classificaram -se com conceito 4, o valor modal.
Esse foi também o conceito modal da região Sudeste ( 41,9 \% ).
O conceito 3 foi o segundo mais frequente em nível nacional ( 29,0 \%, correspondendo a 27 cursos ), sendo o conceito modal das regiões Sul ( 42,1 \% ) e Centro-Oeste ( 50,0 \% ), e o conceito 2, o terceiro ( 20,4 \%, correspondendo a 19 cursos ), foi o conceito modal nas regiões Norte ( 100,0 \% ) e Nordeste ( 35,5 \% ).
Houve, ainda, 14 cursos ( 15,1 \% ) que receberam conceito 5, e dois cursos ( 2,2 \% ) que receberam conceito 1.
Dos 93 cursos de Ciências Sociais, um ( 1,1 \% ) ficou Sem Conceito ( SC ).
A região Norte participou com dois cursos ou 2,2 \% do total nacional, os quais receberam conceito 2.
A região Nordeste participou com 31 cursos ou 33,3 \% do total nacional.
Nessa região, 11 cursos ( 35,5 \% em termos regionais ) obtiveram conceito 2, o conceito modal para a região, como já comentado.
Os conceitos 1 e 3 foram atribuídos, respectivamente, a um e a dez cursos ( 3,2 \% e 32,3 \%, respectivamente ).
Ainda nessa região, oito cursos ( 25,8 \% ) foram avaliados com o conceito 4, e um curso ( 3,2 \% ) foi avaliado com o conceito 5.
Nenhum curso ficou Sem Conceito ( SC ).
Dos 31 cursos participantes da região Sudeste ( 33,3 \% do total nacional ), o conceito 4 foi o conceito modal, como já anteriormente assinalado, tendo sido obtido por 13 cursos ( 41,9 \% ).
O conceito 2 foi atribuído a cinco cursos ( 16,1 \% ).
Receberam os conceitos 3 e 5, respectivamente, quatro cursos ( 12,9 \% ) e nove cursos ( 29,0 \% ).
Nessa região, nenhum curso ficou Sem Conceito ( SC ) ou recebeu conceito 1.
Os 19 cursos da região Sul corresponderam a 20,4 \% do total nacional.
Como já indicado, houve predominância do conceito 3, atribuído a 42,1 \% dos cursos da região ( conceito modal ), ou seja, atribuído a oito dos 19 cursos participantes na região Sul.
O conceito 4 foi atribuído a sete cursos ( 36,8 \% ), e o conceito 5, a quatro cursos ( 21,1 \% ).
Nenhum curso ficou Sem Conceito ( SC ) na região Sul ou recebeu conceito 1 ou 2.
Já dos dez cursos participantes na região Centro-Oeste ( 10,8 \% do total nacional ), cinco ( 50,0 \% em termos regionais ) receberam conceito 3, o conceito modal, como já fora destacado.
Um curso ( 10,0 \% ) obteve o conceito 1, um curso ( 10,0 \% ) obteve conceito 2 e dois cursos obtiveram o conceito 4.
Um curso ficou Sem Conceito ( SC ) e nenhum recebeu o conceito 5.
No Gráfico 5.1a, é apresentada a distribuição cumulativa do Conceito Enade dos cursos de Licenciatura segundo Grande Região.
Em um gráfico de distribuição cumulativa, em cada valor de abscissa ( eixo x ), o valor apresentado no eixo das ordenadas ( eixo y ) é o valor acumulado da variável ( Conceito Enade ).
Observando a linha vertical pontilhada, que \textbf{ajuda} a visualizar o percentual de conceitos acumulados até o valor 3, ou seja, conceitos 1, 2 e 3, pode -se notar que, por exemplo, a região Norte ( linha azul ) apresenta 100 \% dos cursos com conceito abaixo deste valor.
A região Sudeste ( linha vermelha ), por outro lado, apresenta menos de 30 \% dos cursos na mesma situação : com conceito 3 ou abaixo.
Destaca -se que em algumas regiões o conceito 5 pode não ter sido alcançado por nenhum curso.
É o que ocorre na região Norte ( linha azul ), já mencionado, na qual todos os cursos têm conceito igual a 2, e a poligonal alcança 100 \% no conceito 2, e na região Centro-Oeste ( linha \textbf{laranja} ), na qual todos os cursos têm conceito menor ou igual a 4, e a poligonal alcança 100 \% no conceito 4.
Em linhas gerais, em um gráfico de distribuição cumulativa, poligonais mais à esquerda representam Grandes Regiões com pior distribuição, e poligonais mais à direita, Grandes Regiões com melhores conceitos.
A região Norte ( linha azul ) apresenta os cursos com conceitos mais baixos, ao passo que a região Sudeste ( linha vermelha ) apresenta a distribuição com valores maiores.
Cursos em IES situadas nas regiões Nordeste ( linha verde ), Sul ( linha rosa ) e Centro-Oeste ( linha \textbf{laranja} ) apresentam uma situação intermediária.
Na Tabela 5.1b, são apresentadas a quantidade e a distribuição de cursos de Bacharelado de Ciências Sociais participantes do Enade/2021, por faixa de conceito e Grande Região.
Observando -se os dados da Tabela 5.1b, nota -se que, dos 71 cursos participantes, 22 ( 31,0 \% ) classificaram -se com conceito 4, o valor modal.
Esse foi também o conceito modal nas regiões Sudeste ( 34,6 \% ) e Centro-Oeste ( 50,0 \% ).
O conceito 3 foi o segundo mais frequente em nível nacional ( 29,6 \%, correspondendo a 21 cursos ), sendo esse o conceito modal nas regiões Nordeste ( 39,1 \% ) e Sul ( 50,0 \% ).
Já o conceito 2 foi o terceiro mais frequente ( 25,4 \%, correspondendo a 18 cursos ), sendo esse o conceito modal na região Norte ( 83,3 \% ).
Houve, ainda, quatro cursos ( 5,6 \% ) que receberam conceito 1, e seis cursos que receberam o conceito 5.
Nenhum curso de Bacharelado em Ciências Sociais ficou Sem Conceito ( SC ).
A região Norte participou com seis cursos ou 8,5 \% do total nacional.
Desses, cinco cursos receberam o conceito 2, o conceito modal, como já assinalado, o que equivale a 83,3 \% do total regional, e a um curso ( 16,7 \% ) atribuiu -se conceito 1.
Nenhum curso recebeu os conceitos 3, 4 ou 5.
A região Nordeste participou com 23 cursos ou 32,4 \% do total nacional.
Nessa região, nove cursos ( 39,1 \% em termos regionais ) obtiveram conceito 3, o conceito modal para a região, como já comentado.
Os conceitos 1 e 2 foram atribuídos, respectivamente, a um e a sete cursos ( 4,3 \% e 30,4 \%, respectivamente ).
Ainda nessa região, seis cursos ( 26,1 \% ) foram avaliados com o conceito 4.
Nenhum curso recebeu o conceito 5.
Dos 26 cursos participantes da região Sudeste ( 36,6 \% do total nacional ), o conceito 4 foi o conceito modal, como já anteriormente assinalado, tendo sido obtido por nove cursos ( 34,6 \% ).
O conceito 1 foi atribuído a um curso ( 3,8 \% ), e o conceito 2, a cinco cursos ( 19,2 \% ).
Receberam os conceitos 3 e 5, respectivamente, seis cursos ( 23,1 \% ) e cinco cursos ( 19,2 \% ).
Os oito cursos da região Sul corresponderam a 11,3 \% do total nacional.
Como já indicado, houve predominância do conceito 3, atribuído a 50,0 \% dos cursos da região ( conceito modal ), ou seja, atribuído a quatro dos oito cursos participantes na região Sul.
O conceito 4 foi atribuído a três cursos ( 37,5 \% ), e o conceito 5, a um curso ( 12,5 \% ).
Nenhum curso recebeu os conceitos 1 ou 2.
Já dos oito cursos participantes na região Centro-Oeste ( 11,3 \% do total nacional ), quatro ( 50,0 \% em termos regionais ) receberam conceito 4, o conceito modal, como já anteriormente comentado.
Um curso ( 12,5 \% ) obteve conceito 1, um ( 12,5 \% ) obteve conceito 2 e dois cursos ( 25,0 \% ) obtiveram conceito 3.
Nenhum curso recebeu conceito 5.
No Gráfico 5.1b, é apresentada a distribuição cumulativa do Conceito Enade dos cursos de Bacharelado segundo Grande Região.
Em um gráfico de distribuição cumulativa, em cada valor de abscissa ( eixo x ), o valor apresentado no eixo das ordenadas ( eixo y ) é o valor acumulado da variável ( conceito Enade ).
Observando a linha vertical pontilhada, que \textbf{ajuda} a visualizar o percentual de conceitos acumulados até o valor 3, ou seja, conceitos 1, 2 e 3, pode -se notar que, por exemplo, a região Norte ( linha azul ) apresenta 100 \% dos cursos com conceito abaixo deste valor.
Já região Centro-Oeste ( linha \textbf{laranja} ), por outro lado, apresenta 50 \% dos cursos : com conceito 3 ou abaixo.
Destaca -se que em algumas regiões o conceito 5 pode não ter sido alcançado por nenhum curso, como, já comentado, ocorre na região Norte ( linha azul ), na qual todos os cursos têm conceito menor ou igual a 2, e sua poligonal alcança 100 \% neste conceito, e nas regiões Nordeste e Centro-Oeste, nas quais todos os cursos têm conceito menor ou igual a 4, e sua poligonal alcança 100 \% neste conceito.
Em linhas gerais, em um gráfico de distribuição cumulativa, poligonais mais à esquerda representam Grandes Regiões com pior distribuição, e poligonais mais à direita, Grandes Regiões com melhores conceitos.
A região Norte ( linha azul ) apresenta os cursos com conceitos mais baixos, ao passo que as regiões Sudeste ( linha vermelha ) e Sul ( linha rosa ) apresentam as distribuições com valores maiores.
Cursos em IES situadas nas regiões Nordeste ( linha verde ) e Centro-Oeste ( linha \textbf{laranja} ) apresentam uma situação intermediária. 5.2 CONCEITOS POR CATEGORIA ADMINISTRATIVA, POR MODALIDADE DE ENSINO E GRANDE REGIÃO Na Tabela 5.2a, é apresentada a distribuição dos cursos dos cursos de Licenciatura de Ciências Sociais participantes do Enade/2021, por Categoria Administrativa e Modalidade de Ensino, de acordo com os conceitos por eles alcançados, segundo as Grandes Regiões brasileiras.
Dos 93 cursos participantes, 74 ( 79,6 \% ) eram ministrados em Instituições Públicas, e 19 ( 20,4 \% ), em Instituições Privadas.
Quanto à Modalidade de Ensino, 77 ( 82,8 \% ) cursos eram Presenciais e 16 ( 17,2 \% ), a Distância.
De acordo com as informações da Tabela 5.2a, em termos nacionais, entre Instituições Públicas e Privadas, observa -se que, dos 14 cursos avaliados com conceito 5, 12 eram oferecidos em IES Públicas, e dois, em IES Privadas.
Dos 74 cursos participantes de IES Públicas, o conceito 4 foi o valor modal, atribuído a 23 cursos ( 31,1 \% ).
Dos demais cursos de Instituições Públicas participantes, dois ( 2,7 \% ) foram avaliados com conceito 1, 15 cursos ( 20,3 \% ) com conceito 2, 21 cursos ( 28,4 \% ) com conceito 3 e, como já comentado, 12 cursos ( 16,2 \% ), com conceito 5.
Nesta categoria, um curso ficou Sem Conceito ( SC ).
Na rede privada, o conceito modal foi 4, atribuído a sete cursos, o correspondente a 36,8 \% dos 19 cursos da categoria.
Dos demais cursos participantes, quatro ( 21,1 \% ) receberam conceito 2, seis ( 31,6 \% ), conceito 3 e dois ( 10,5 \% ), conceito 5.
Nesta Categoria Administrativa, nenhum curso ficou Sem Conceito ( SC ) ou recebeu conceito 1.
Dos 77 cursos participantes da Modalidade Presencial, o conceito 4 foi o valor modal, atribuído a 27 cursos ( 35,1 \% ).
Dos demais cursos nesta Modalidade de Ensino, um curso ( 1,3 \% ) recebeu conceito 1, 14 cursos ( 18,2 \% ) receberam conceito 2, 20 cursos ( 26,0 \% ), conceito 3, e 14 cursos ( 18,2 \% ), conceito 5.
Nesta Modalidade de Ensino, um curso ( 1,4 \% ) ficou Sem Conceito ( SC ).
Foram 16 cursos a Distância informados, e a sete ( 43,8 \% ) atribuiu -se o conceito modal 3.
Nessa Modalidade de Ensino, nenhum curso ficou Sem Conceito ( SC ) ou recebeu conceito 5.
Dos demais cursos nesta Modalidade de Ensino, um curso ( 6,3 \% ) recebeu conceito 1, cinco cursos ( 31,3 \% ) receberam conceito 2, e três cursos ( 18,8 \% ), conceito 4.
Na análise por região, observa -se que, na região Norte, não houve cursos em Instituições Privadas, e os dois cursos foram de Instituições Públicas, na Modalidade Educação Presencial, avaliados com conceito 2.
Na região Nordeste, também todos os 31 cursos foram de Instituições Públicas e 11 cursos obtiveram o conceito modal 2.
Um curso obteve conceito 1, dez receberam o conceito 3, oito cursos receberam conceito 4 e um curso ficou com conceito 5, não havendo nenhum curso da região Nordeste que tenha ficado Sem Conceito ( SC ).
Considerando a Modalidade de Ensino, observa -se que, na região Nordeste, dos 31 cursos oferecidos, 30 ( 96,8 \% do total regional ) eram Presenciais, e, destes, 11 ( 36,7 \% ) obtiveram conceito 2 ( valor modal ).
Dez cursos ( 33,3 \% ) receberam conceito 3, oito cursos ( 26,7 \% ) receberam conceito 4 e um curso ( 3,3 \% ) recebeu conceito 5, não tendo nenhum curso ficado Sem Conceito ( SC ) ou recebido conceito 1.
Já o único curso ( 3,2 \% do total regional ) oferecido nessa região na Modalidade a Distância obteve conceito 1.
Na região Sudeste, a proporção de cursos da rede privada foi 35,5 \%, correspondendo a 11 dos 31 cursos participantes.
Nessa categoria e região, o conceito modal foi 4, atribuído a cinco cursos.
Os demais foram avaliados com conceito 2 ( quatro cursos ) e conceito 3 ( dois cursos ).
Nenhum curso de Instituição Privada nessa região ficou Sem Conceito ( SC ) ou obteve conceito 1 ou 5.
Dos 20 cursos oferecidos por Instituições Públicas na região Sudeste, o conceito modal foi 5, atribuído a nove cursos.
Nessa região, os demais conceitos atribuídos foram o conceito 2 ( um curso ), o conceito 3 ( dois cursos ) e o conceito 4 ( oito cursos ).
Nenhum curso de Instituição Pública nessa região ficou Sem Conceito ( SC ) ou obteve conceito 1.
Também na região Sudeste, 22 dos 31 cursos participantes foram oferecidos na Modalidade Presencial.
Para esses, a moda foi o conceito 4, com 11 cursos.
Os demais cursos Presenciais da região foram alocados nos conceitos 3 ( dois cursos ) e 5 ( nove cursos ), não tendo nenhum curso ficado Sem Conceito ( SC ) ou recebido os conceitos 1 ou 2.
Já dos nove cursos participantes na Modalidade Educação a Distância na região Sudeste, cinco obtiveram conceito 2, e os conceitos 3 e 4 foram atribuídos a dois cursos cada, não havendo cursos Sem Conceito ( SC ) ou com os conceitos 1 ou 5.
As Instituições Privadas concentraram sete dos 19 cursos participantes da região Sul, 36,8 \% do total regional.
Desses, três cursos obtiveram conceito 3, o conceito modal.
Nesta combinação de Categoria Administrativa e Grande Região, dois cursos receberam conceito 4 e outros dois, conceito 5.
Nenhum curso recebeu os conceitos 1 ou 2 ou ficou Sem Conceito ( SC ).
As Instituições Públicas na região Sul participaram com 12 cursos ( 63,2 \% dos cursos da região ), aos quais foram atribuídos os conceitos modais 3 e 4 a cinco cursos ( distribuição bimodal ).
Os dois cursos restantes obtiveram conceito 5.
Nenhum curso da região Sul oferecido por IES Públicas ficou Sem Conceito ( SC ) ou recebeu os conceitos 1 e 2.
Quanto à Modalidade de Ensino, também na região Sul, 14 dos 19 cursos participantes foram oferecidos na Modalidade Presencial.
Para esses, a moda foi o conceito 4, com seis cursos.
Os demais cursos Presenciais da região foram alocados nos conceitos 3 ( quatro cursos ) e 5 ( quatro cursos ).
Nenhum curso ficou Sem Conceito ( SC ) ou recebeu os conceitos 1 ou 2.
Já dos cinco cursos participantes na Modalidade Educação a Distância na região Sul, quatro obtiveram conceito 3 ( conceito modal ) e um obteve conceito 4.
Na região Centro Oeste, apenas um dos dez cursos participantes eram de instituições Privadas ( 10,0 \% em termos regionais ), e recebeu conceito 3.
Quanto aos nove cursos oferecidos por Instituições Públicas na região Centro Oeste ( 90,0 \% do total regional ), quatro apresentaram conceito 3, conceito modal.
Nesta combinação de Categoria Administrativa e Grande Região, dois cursos receberam conceito 4, os conceitos 1 e 2 foram atribuídos a um curso cada e um curso ficou Sem Conceito ( SC ), e nenhum curso recebeu o conceito 5.
No que se refere à Modalidade de Ensino, nove dos 10 cursos participantes eram presenciais.
Para esses, a moda foi o conceito 3, com quatro cursos.
Os demais cursos presenciais da região foram alocados nos conceitos 1 ( um curso ), 2 ( um curso ) e 4 ( dois cursos ).
Um curso ficou Sem Conceito ( SC ), não tendo nenhum curso recebido o conceito 5.
Já o único curso participante na Modalidade Educação a Distância na região Centro-Oeste, obteve conceito 3.
No Gráfico 5.2a, apresenta -se a distribuição do Conceito Enade dos cursos de Licenciatura, segundo a Categoria Administrativa da IES.
Os cursos em IES Públicas ( linha azul ) apresentam uma poligonal mais à direita do que os em IES Privadas ( linha verde ) e, consequentemente, uma distribuição de Conceito Enade com valores maiores.
Na Tabela 5.2b, é apresentada a distribuição dos cursos dos cursos de Bacharelado de Ciências Sociais participantes do Enade/2021, por Categoria Administrativa e Modalidade de Ensino, de acordo com os conceitos por eles alcançados, segundo as Grandes Regiões brasileiras.
Dos 71 cursos participantes, 63 ( 88,7 \% ) eram ministrados em Instituições Públicas, e oito ( 11,3 \% ), em Instituições Privadas.
Quanto à Modalidade de Ensino, todos os 71 cursos oferecidos eram Presenciais.
De acordo com as informações da Tabela 5.2b, em termos nacionais, entre Instituições Públicas e Privadas, observa -se que, dos seis cursos avaliados com conceito 5, quatro eram oferecidos em IES Públicas, e dois, em IES Privadas.
Dos 63 cursos participantes de IES Públicas, o conceito 4 foi o valor modal, atribuído a 21 cursos ( 33,3 \% ).
Dos demais cursos de Instituições Públicas participantes, quatro cursos ( 6,3 \% ) foram avaliados com conceito 1, 16 cursos ( 25,4 \% ) foram avaliados com conceito 2, 18 cursos ( 28,6 \% ) foram avaliados com conceito 3 e, quatro cursos ( 6,3 \% ) receberam conceito 5.
Na rede privada, o conceito modal foi 3, atribuído a três cursos, o correspondente a 37,5 \% dos oito cursos da categoria.
Dos demais cursos participantes, dois ( 25,0 \% ) receberam conceito 2, um curso ( 12,5 \% ) foi avaliado com conceito 4 e dois cursos ( 25,0 \% ) receberam conceito 5.
Nesta Categoria Administrativa, nenhum curso ficou Sem Conceito ( SC ) ou recebeu conceito 1.
Conforme mencionado, todos os 71 cursos participantes eram da Modalidade Presencial, e o conceito 4 foi o valor modal, atribuído a 22 cursos ( 31,0 \% ).
Dos demais cursos nesta Modalidade de Ensino, quatro cursos ( 5,6 \% ) receberam conceito 1, 18 cursos ( 25,4 \% ) receberam conceito 2, 21 cursos ( 29,6 \% ), conceito 3, e seis cursos ( 8,5 \% ), conceito 5.
Nesta Modalidade de Ensino, nenhum curso ficou Sem Conceito ( SC ).
Na análise por região, observa -se que, na região Norte, todos os seis cursos participantes eram vinculados a Instituições Públicas e eram Presenciais.
O conceito modal foi 2, atribuído a cinco cursos, e um curso obteve conceito 1.
Na região Nordeste, apenas um dos 23 cursos participantes era rede privada, e recebeu conceito 2.
As Instituições Públicas da região Nordeste participaram com 22 cursos ( 95,7 \% do total da região ).
Desses, nove obtiveram conceito 3, o valor modal, e os conceitos 2 e 4 foram atribuídos a seis cursos cada.
Um curso recebeu conceito 1.
Nenhum curso recebeu conceito 5.
Na região Sudeste, a proporção de cursos da rede privada, 23,1 \%, corresponde a seis dos 26 cursos participantes.
Nessa categoria e região, o conceito modal foi 3, atribuído a três cursos.
Os demais foram avaliados com conceito 2 ( um curso ) e conceito 5 ( dois cursos ).
Nenhum curso foi avaliado com conceitos 1 ou 4.
Dos 20 cursos oferecidos por Instituições Públicas na região Sudeste, o conceito modal foi 4, atribuído a nove cursos.
Nessa região, um curso de IES Públicas recebeu conceito 1, quatro cursos receberam conceito 2, e os conceitos 3 e 5 foram atribuídos a três cursos cada.
Apenas um dos oito cursos participantes da região Sul eram de Instituições Privadas participaram ( 12,5 \% do total regional ), e obteve conceito 4.
Já as Instituições Públicas na região Sul participaram com sete cursos ( 87,5 \% dos cursos da região ), e quatro cursos receberam o conceito modal 3, dois cursos receberam conceito 4 e um recebeu conceito 5.
Nenhum curso recebeu os demais conceitos.
Na região Centro-Oeste, todos os cursos participantes eram de Instituições Públicas.
Nessa região, quatro cursos apresentaram conceito 4, conceito modal, dois receberam conceito 3, e os conceitos 1 e 2 foram atribuídos a um curso cada.
Nenhum curso recebeu conceito 5.
No Gráfico 5.2b, apresenta -se a distribuição do Conceito Enade dos cursos de Bacharelado, segundo a Categoria Administrativa da IES.
Os cursos em IES Privadas ( linha verde ) apresentam uma poligonal mais à direita do que os em IES Públicas ( linha azul ) e, consequentemente, uma distribuição de Conceitos Enade com valores maiores. 5.3 CONCEITOS POR ORGANIZAÇÃO ACADÊMICA E GRANDE REGIÃO Na Tabela 5.3a, encontra -se a distribuição dos conceitos atribuídos aos cursos de Licenciatura participantes do Enade/2021, na Área de Ciências Sociais, por Organização Acadêmica, segundo as Grandes Regiões brasileiras.
Dos 93 cursos de Ciências Sociais participantes, 86 eram oferecidos em Universidades, quatro em Centros Universitários, três em CEFET/IF.
Esta distribuição corresponde a, respectivamente, 92,5 \%, 4,3 \% e 3,2 \% dos cursos.
Nenhum curso foi oferecido em Faculdades.
De acordo com os dados apresentados, todos os 14 cursos avaliados com conceito 5 eram vinculados a Universidades.
Os cursos vinculados a esse tipo de Organização Acadêmica tiveram o conceito 4 como conceito modal, atribuído a 29 cursos ( 33,7 \% ).
Os demais cursos vinculados a Universidades receberam os conceitos 1 ( dois cursos ), 2 ( 18 cursos ), 3 ( 22 cursos ) e conceito 5 ( 14 cursos, como já mencionado ).
Um curso ficou Sem Conceito ( SC ).
Dos quatro cursos em Centros Universitários, o conceito modal foi 3, atribuído a dois cursos ( 50,0 \% ).
Nenhum curso vinculado a esse tipo de Organização Acadêmica ficou Sem Conceito ( SC ), e os outros dois cursos receberam conceitos 2 e 4.
Todos os três cursos mantidos por CEFET/IF receberam conceito 3.
Como já foi mencionado, não houve cursos de Licenciatura na Área de Ciências Sociais avaliados em Faculdades no Enade/2021.
Considerando -se, separadamente, as regiões brasileiras, verifica -se que, na região Norte, os dois cursos participantes eram vinculados a Universidades, e a eles foi atribuído conceito 2.
Na região Nordeste, também todos os cursos participantes ( 31 ) eram vinculados a Universidades, e nenhum ficou Sem Conceito ( SC ).
O conceito modal 2 foi atribuído a 11 cursos.
Um curso obteve conceito 1, dez cursos, conceito 3, oito cursos, conceito 4, e um curso obteve conceito 5.
Na região Sudeste, as Universidades concentraram 29 dos 31 cursos da região.
Dos cursos oferecidos em Universidades, nessa região, o conceito modal foi 4, atribuído a 12 cursos.
Os demais cursos receberam os conceitos 2 ( quatro cursos ), 3 ( quatro cursos ) e 5 ( nove cursos ).
Nenhum curso ficou Sem Conceito ( SC ) ou recebeu conceito 1.
Os Centros Universitários participaram com dois cursos restantes na região Sudeste, dos quais um obteve o conceito 2 e, o outro, conceito 4.
Dos 19 cursos da região Sul, 16 eram mantidos por Universidades, para os quais o conceito modal foi 4, atribuído a sete cursos.
Os demais cursos receberam os conceitos 3 ( cinco cursos ) e 5 ( quatro cursos ).
Nesse tipo de organização, nenhum curso da região Sul ficou Sem Conceito ( SC ) ou recebeu conceitos 1 ou 2.
Os dois cursos participantes de Centros Universitários da região Sul receberam o conceito 3, mesmo conceito obtido pelo único curso participante vinculado a CEFET/IF nessa região.
Na região Centro-Oeste, oito dos dez cursos oferecidos eram mantidos por Universidades.
Para este tipo de Organização Acadêmica, o conceito modal foi 3, atribuído a três cursos, e os demais receberam conceito 1 ( um curso ), conceito 2 ( um curso ) e conceito 4 ( dois cursos ).
Nesse tipo de organização na região Centro-Oeste, um curso ficou Sem Conceito ( SC ) e nenhum obteve conceito 5.
Já os dois cursos restantes nessa região eram vinculados a CEFET/IF e receberam conceito 3.
No Gráfico 5.3a, apresenta -se a distribuição do Conceito Enade dos cursos de Licenciatura, segundo a Organização Acadêmica da IES.
Os cursos em Universidades ( linha azul ) apresentam uma poligonal mais à direita do que as demais, denotando uma melhor distribuição de Conceitos Enade.
No outro extremo, encontram -se os cursos mantidos por CEFET/IF ( linha vermelha ), apresentando as respectivas poligonais mais à esquerda e, consequentemente, distribuição de Conceitos Enade com menores valores.
Os cursos mantidos por Centros Universitários ( linha verde ) apresentam poligonais mais centrais e, consequentemente, distribuição de Conceitos Enade com valores intermediários.
Na Tabela 5.3b, encontra -se a distribuição dos conceitos atribuídos aos cursos de Bacharelado participantes do Enade/2021, na Área de Ciências Sociais, por Organização Acadêmica, segundo as Grandes Regiões brasileiras.
Dos 71 cursos de Ciências Sociais participantes, 68 eram oferecidos em Universidades, um em Centros Universitários e dois em Faculdades.
Esta distribuição corresponde a, respectivamente, 95,8 \%, 1,4 \% e 2,8 \% dos cursos.
Nenhum curso foi oferecido em CEFET/IF.
De acordo com os dados apresentados, de todos os seis cursos avaliados com conceito 5, cinco eram vinculados a Universidades.
Os cursos vinculados a esse tipo de Organização Acadêmica tiveram o conceito 4 como conceito modal, atribuído a 22 cursos ( 32,4 \% ).
Os demais cursos vinculados a Universidades receberam os conceitos 1 ( quatro cursos ), 2 ( 16 cursos ), 3 ( 21 cursos ) e conceito 5 ( cinco cursos, como já mencionado ).
Nenhum curso ficou Sem Conceito ( SC ).
Apenas um curso era vinculado a Centros Universitários, e obteve o conceito 2.
Dos dois cursos mantidos por Faculdades, um recebeu conceito 2 e, o outro, conceito 5.
Considerando -se, separadamente, as regiões brasileiras, verifica -se que, na região Norte, todos os seis cursos participantes eram mantidos por Universidades.
O conceito modal foi 2, atribuído a cinco cursos, e um curso recebeu conceito 1.
Na região Nordeste, as Universidades participaram com 22 dos 23 cursos da Área de Ciências Sociais ( Bacharelado ) oferecidos.
O conceito modal foi 3, com nove cursos.
O conceito 1 foi atribuído a um curso, e os conceitos 2 e 4, a seis cursos cada.
Nenhum curso recebeu conceito 5.
Os Centros Universitários contaram com um curso participante na região Nordeste, o qual recebeu o conceito 2.
Nessa região não houve cursos avaliados mantidos por Faculdades.
Na região Sudeste, as Universidades concentraram 24 dos 26 cursos de Ciências Sociais ( Bacharelado ) da região.
Dos cursos oferecidos em Universidades, nessa região, o conceito modal foi 4, atribuído a nove cursos.
Os demais cursos receberam os conceitos 1 ( um curso ), 2 ( quatro cursos ), 3 ( seis cursos ) e 5 ( quatro cursos ).
Já as Faculdades foram representadas por dois cursos na região Sudeste, dos quais um recebeu o conceito 2 e, o outro, conceito 5.
Nessa região não houve cursos avaliados mantidos por Centros Universitários.
Todos os oito cursos da região Sul eram mantidos por Universidades, para os quais o conceito modal foi 3, atribuído a quatro cursos.
O segundo conceito mais frequente foi 4, com três cursos, e o conceito 5 foi obtido por um curso.
Nenhum curso recebeu os conceitos 1 ou 2.
Na região Centro-Oeste, também os 8 cursos avaliados eram mantidos por Universidades.
O conceito modal foi 4 ( quatro cursos ), e os demais cursos receberam os conceitos 1 ( um curso ), 2 ( um curso ) e 3 ( dois cursos ).
Nenhum curso recebeu conceito 5.
No Gráfico 5.3b, apresenta -se a distribuição do Conceito Enade dos cursos de Bacharelado, segundo a Organização Acadêmica da IES.
Nos Centros Universitários ( linha verde ), todos os cursos têm conceito 2 ou abaixo disso, já que a poligonal cumulativa alcança 100 \% para o conceito 2.
Os cursos em Faculdades ( linha vermelha ) apresentam uma poligonal mais à direita do que as demais, denotando uma melhor distribuição de Conceitos Enade.
No outro extremo encontram -se os cursos mantidos por Centros Universitários ( linha verde ), apresentando poligonal mais à esquerda e, consequentemente, distribuição de Conceitos Enade com menores valores.
Os cursos mantidos por Universidades ( em azul ) apresentam poligonais mais centrais e, consequentemente, distribuição de Conceitos Enade com valores intermediários.

% matched lemmas: ajudar, laranja
\end{textsample}
